\section{A counterexample to the Auslander--Reiten conjecture}
\label{sec:ar-route}

Following~\cite{OAI26ar}, we construct a bimodule and a stable morphism
satisfying the hypotheses of \Cref{thm:conversion}, and deduce a second algebra
of infinite little finitistic dimension, with nine simple modules. Throughout,
\[
  \kk=\mathbb{F}_2(q,H_1,H_2),
\]
the field of rational functions in three variables over $\mathbb{F}_2$; in
characteristic two all signs disappear, and ${\delta^a(g,j)=F(g)v+v[a]j}$ in
\Cref{thm:conversion}. The construction uses a ten-dimensional algebra $C$ with
a simple module $s$ whose resolution has differentials depending on successive
powers of $q$ (\Cref{subsec:algebra-C}); its trivial extension ${T=C\ltimes DC}$,
over which $s$ has polynomial Ext algebra on a generator of degree three
(\Cref{subsec:algebra-T}), realised by an explicit Hochschild cocycle
(\Cref{subsec:cocycle}); the tensor square ${E=T\otimes_\kk T}$, a bimodule
$\mathcal{C}$ built from two cones on that cocycle, and twisted lifts
(\Cref{subsec:tensor-square,subsec:two-cones,subsec:lifts}); and the fibre $F$
of a map from two twisted copies of $E$ to $\mathcal{C}[3]$
(\Cref{subsec:fibre}). Finite identities, such as associativity of $C$ and
the cocycle identities, are checked by exact computations recorded in
\Cref{app:computations}; every statement about all degrees is proved in the
text.

\subsection{The algebra \texorpdfstring{$C$}{C}}
\label{subsec:algebra-C}

Let $C$ be the $\kk$-algebra with basis
\[
  e,x,y,z\in eCe,\qquad u,v\in eCf,\qquad t,j\in fCe,\qquad f,n\in fCf,
\]
where $e$ and $f$ are orthogonal idempotents with ${e+f=1}$, multiplication by
$e$ and $f$ is prescribed by the corners, and the products of two of the
remaining eight basis vectors are zero except for
\[
  xy=qz,\quad yx=z,\quad xu=yu=v,\quad tx=j,\quad ty=q^2j,\quad ut=y+qx,
\]
\[
  vt=qz,\quad uj=z,\quad tu=(1+q)n,\quad nt=qj,\quad un=v.
\]
Let $N$ be the span of $x,y,z,u,v,t,j,n$, let $s$ be the simple left module
${C/(N+Ce)}$ at $f$, and let $s^{\mathrm{r}}$ be the simple right module at~$f$.

\begin{lemma}
  \label{lemma:algebra-C}
  The algebra $C$ is associative. It is graded by
  ${\deg u=\deg t=1}$, ${\deg x=\deg y=\deg n=2}$, ${\deg v=\deg j=3}$, ${\deg z=4}$
  and ${\deg e=\deg f=0}$, its radical is $N$, and ${N^5=0}$.
\end{lemma}

\begin{proof}
  Every product above is homogeneous for the grading. A product of three
  basis vectors vanishes under both bracketings if one of them is an
  idempotent that does not match the corners, or if its total degree exceeds
  four, since there is no basis vector of degree larger than four. Among the
  remaining triples of radical basis vectors, the only ones with compatible
  corners are $xut$, $yut$, $utx$, $uty$, $utu$, $unt$, $txu$, $tyu$, $tut$,
  $tun$ and $ntu$, and their products under both bracketings are $qz$, $qz$,
  $z$, $q^2z$, ${(1+q)v}$, $qz$, $0$, $0$, ${q(1+q)j}$, $0$ and $0$; for instance
  ${(ut)y=(y+qx)y=q^2z=u(q^2j)=u(ty)}$. The associators of all $1000$ triples
  of basis vectors are also checked by an exact computation
  (\Cref{app:computations}). The span $N$ of the basis vectors of positive
  degree is a nilpotent ideal with ${N^5=0}$ and ${C/N\cong\kk\times\kk}$, so $N$
  is the radical.
\end{proof}

For ${i\geq0}$ put ${\ell_i=x+q^iy\in eCe}$. The left modules $Ce$ and $Cf$ have
the bases $e,x,y,z,t,j$ and $f,u,v,n$, and right multiplication by $u$ and by
$\ell_i$ acts as follows:
\begin{equation}
  \label{eq:C-right-multiplication}
  \begin{array}{c|cccccc}
    a&e&x&y&z&t&j\\\hline
    au&u&v&v&0&(1+q)n&0\\
    a\ell_i&\ell_i&q^{i+1}z&z&0&(1+q^{i+2})j&0
  \end{array}
\end{equation}

\begin{proposition}
  \label{prop:C-resolution}
  The complex
  \[
    \mathbb{R}\colon\quad\cdots\longrightarrow Ce\xrightarrow{\cdot\ell_2}Ce
    \xrightarrow{\cdot\ell_1}Ce\xrightarrow{\cdot\ell_0}Ce
    \xrightarrow{\cdot u}Cf,
  \]
  with $Cf$ in degree zero and right multiplications as differentials, is a
  minimal projective resolution of $s$. Consequently,
  ${\operatorname{Ext}^a_C(s,s)=0}$ for ${a>0}$, and ${\RHom[C]{s}{C}\cong
  s^{\mathrm{r}}[-2]}$ in the derived category of right $C$-modules.
\end{proposition}

\begin{proof}
  Since $q$ is transcendental over $\mathbb{F}_2$, the elements ${1+q}$ and
  ${1+q^{i+2}}$ are nonzero. By~\eqref{eq:C-right-multiplication}, the image of
  ${\cdot u}$ is the span of $u$, $v$ and $n$, which is the radical of $Cf$, and
  its kernel is spanned by $\ell_0$, $z$ and $j$. A vector in the kernel of
  ${\cdot\ell_i}$ has zero coefficients at $e$ and $t$, and coefficients at $x$
  and $y$ related by ${a_y=q^{i+1}a_x}$; hence
  \[
    \ker(\cdot\ell_i)=\langle\ell_{i+1},z,j\rangle,\qquad
    \operatorname{im}(\cdot\ell_i)=\langle\ell_i,z,j\rangle,
  \]
  for every ${i\geq0}$, and the complex is exact. All differentials have image
  in $N$ times their targets, so each surjection onto a syzygy is a projective
  cover, and the resolution is minimal. Since ${\Hom[C]{Ce}{s}=es=0}$ and
  ${\Hom[C]{Cf}{s}=\kk}$, the asserted vanishing of self-extensions follows.

  Identifying ${\Hom[C]{Ce}{C}=eC}$ and ${\Hom[C]{Cf}{C}=fC}$ as right modules,
  precomposition with right multiplication by $b$ becomes left multiplication
  by $b$, and $\Hom[C]{\mathbb{R}}{C}$ is the complex
  \[
    fC\xrightarrow{u\cdot}eC\xrightarrow{\ell_0\cdot}eC\xrightarrow{\ell_1\cdot}
    eC\longrightarrow\cdots
  \]
  in degrees $0,1,2,\dots$, with ${fC=\langle f,t,j,n\rangle}$ and
  ${eC=\langle e,x,y,z,u,v\rangle}$. Left multiplication by $u$ sends
  $f,t,j,n$ to the independent vectors ${u,y+qx,z,v}$, and left multiplication
  by $\ell_i$ sends $e,x,y,z,u,v$ to ${\ell_i,q^iz,qz,0,(1+q^i)v,0}$. For ${i=0}$
  its kernel is ${\langle qx+y,z,u,v\rangle}$, the image of ${u\cdot}$, and its
  image is $\langle\ell_0,z\rangle$; for ${i\geq1}$ its kernel is
  $\langle\ell_{i-1},z,v\rangle$ and its image is $\langle\ell_i,z,v\rangle$.
  Hence, the only cohomology is in degree two, spanned by the class of $v$,
  on which $f$ acts as the identity and the radical acts as zero, since
  ${vt=qz}$ is a boundary. Good truncations identify the complex with
  ${s^{\mathrm{r}}[-2]}$.
\end{proof}

The parameter $q$ makes the resolution aperiodic, while the self-extensions of
$s$ vanish in positive degrees: on $e\langle\ell_i,z,j\rangle$ the elements $x$
and $y$ act with ${x=q^{i+1}y}$ and ${y\neq0}$, a ratio preserved by isomorphisms of
$C$-modules, so the syzygies $\langle\ell_i,z,j\rangle$ are pairwise
nonisomorphic. The module $s$ is not a counterexample to the Auslander--Reiten
conjecture over $C$, since ${\operatorname{Ext}^2_C(s,C)\neq0}$; over ${T=C\ltimes
DC}$, its Ext algebra is polynomial on a generator of degree three
(\Cref{thm:T-ext}).

\subsection{The trivial extension \texorpdfstring{$T$}{T}}
\label{subsec:algebra-T}

Let ${T=C\ltimes DC}$, where the bimodule structure of $DC$ is
${(a\phi)(c)=\phi(ca)}$ and ${(\phi a)(c)=\phi(ac)}$, and regard $s$ as a
$T$-module through ${T\to C}$. For a basis vector $c$ of $C$ we write $c^*$ for
the dual basis vector of $DC$.

\begin{lemma}
  \label{lemma:T-symmetric}
  The algebra $T$ is a twenty-dimensional symmetric algebra, with symmetrising
  form ${\operatorname{tr}(a,\phi)=\phi(1)}$. Its radical is ${N\oplus DC}$, with
  sixth power zero, and ${T/\operatorname{rad}T\cong\kk\times\kk}$.
\end{lemma}

\begin{proof}
  The algebra $T$ is symmetric, with the stated form, by
  \Cref{subsec:trivial-extensions}. With ${\deg c^*=5-\deg c}$, the algebra $T$ is
  graded, with all basis vectors
  other than $e$ and $f$ in degrees between one and five; their span is a
  nilpotent ideal with sixth power zero and quotient ${\kk\times\kk}$.
\end{proof}

\begin{theorem}
  \label{thm:T-ext}
  There is ${\tau\in\operatorname{Ext}^3_T(s,s)}$ such that
  ${\operatorname{Ext}^*_T(s,s)=\kk[\tau]}$, a polynomial algebra under the
  Yoneda product. Every nonzero element of $\operatorname{Ext}^3_T(s,s)$ is
  such a generator. Moreover, $s$ is not projective and
  ${\underline{\operatorname{End}}_T(s)=\kk}$.
\end{theorem}

\begin{proof}
  Let $\mathbb{R}$ be the resolution of \Cref{prop:C-resolution}. For a
  finitely generated projective left $C$-module $V$, the map
  ${DC\otimes_CV\to D\Hom[C]{V}{C}}$, ${\phi\otimes v\mapsto(g\mapsto\phi(g(v)))}$,
  is a natural isomorphism of left $C$-modules, as one sees for ${V=C}$. Hence,
  ${DC\otimes_C\mathbb{R}\cong D\Hom[C]{\mathbb{R}}{C}}$, which has cohomology
  ${D(s^{\mathrm{r}})=s}$ in degree ${-2}$ by \Cref{prop:C-resolution}; since $DC$
  acts on it as zero, it is quasi-isomorphic to $s[2]$ as a complex of
  $T$-modules. Tensoring the sequence ${0\to DC\to T\to C\to0}$ of
  $(T,C)$-bimodules with the complex $\mathbb{R}$ of projective, hence flat,
  $C$-modules gives a short exact sequence of complexes of $T$-modules, and
  hence a triangle
  \begin{equation}
    \label{eq:T-triangle}
    s[2]\longrightarrow\mathcal{I}\longrightarrow s\xrightarrow{\ \eta\ }s[3],
    \qquad\mathcal{I}=T\otimes_C\mathbb{R},
  \end{equation}
  in ${\DerCat[-]{\operatorname{mod}T}}$. The terms ${T\otimes_CCe=Te}$ and
  ${T\otimes_CCf=Tf}$ of $\mathcal{I}$ are projective, and $\mathcal{I}$ is
  bounded above, so adjunction gives
  ${\Hom{\mathcal{I}}{s[a]}\cong H^a\Hom[C]{\mathbb{R}}{s}}$, which is $\kk$ for
  ${a=0}$ and zero otherwise.

  Put ${V^a=\Hom{s}{s[a]}}$, which is $\operatorname{Ext}^a_T(s,s)$ for ${a\geq0}$
  and zero for ${a<0}$. Applying $\Hom{-}{s[a]}$ to~\eqref{eq:T-triangle} gives
  exact sequences
  \[
    \Hom{\mathcal{I}}{s[a-1]}\longrightarrow V^{a-3}
    \xrightarrow{\ \beta\mapsto\beta[3]\circ\eta\ }V^a\longrightarrow
    \Hom{\mathcal{I}}{s[a]}.
  \]
  For ${a=1}$ this gives ${V^1=0}$, and for ${a\geq2}$ the outer terms vanish, so
  that composition with $\eta$ is an isomorphism ${V^{a-3}\to V^a}$. Since
  ${V^0=\kk}$ and ${V^{-1}=V^{-2}=0}$, we obtain ${V^{3m}=\kk\eta^m}$ for ${m\geq0}$
  and ${V^a=0}$ in all other degrees, and no power of $\eta$ vanishes. Hence,
  ${\operatorname{Ext}^*_T(s,s)=\kk[\eta]}$, and every nonzero multiple of $\eta$
  is a generator. Since ${\operatorname{Ext}^3_T(s,s)\neq0}$, the module $s$ is
  not projective, and a nonzero endomorphism of the one-dimensional module $s$
  factoring through a projective module would make $s$ a direct summand of a
  projective module; hence ${\underline{\operatorname{End}}_T(s)=\kk}$.
\end{proof}

\subsection{A Hochschild cocycle}
\label{subsec:cocycle}

The construction of the bimodule $F$ needs a representative of $\tau$ by a
bimodule map, which the abstract computation of \Cref{thm:T-ext} does not
provide. Let ${I=\kk e\oplus\kk f\subseteq T}$ and
${\mathfrak{r}=\operatorname{rad}
T}$, so that ${T=I\oplus\mathfrak{r}}$. The \emph{bar resolution} of $T$ relative
to $I$ is the complex $\mathbb{B}$ of $T$-bimodules with
\[
  \mathbb{B}^{-n}=T\otimes_I\mathfrak{r}^{\otimes_In}\otimes_IT\quad(n\geq0),
  \qquad\mathbb{B}^n=0\quad(n>0),
\]
whose elements we write $a_0[a_1|\cdots|a_n]a_{n+1}$, with the differential
given by the sum of the ${n+1}$ products of adjacent entries, and augmentation
${\mathbb{B}^0\to T}$ given by multiplication. Its terms are finite direct sums
of the projective bimodules ${Ti\otimes_\kk jT}$ for ${i,j\in\set{e,f}}$. The map
\[
  a_0[a_1|\cdots|a_n]a_{n+1}\longmapsto[\overline{a}_0|a_1|\cdots|a_n]a_{n+1},
\]
where $\overline{a}_0$ is the component of $a_0$ in $\mathfrak{r}$, together
with ${a\mapsto1[\,]a}$ in augmented degree, is a contraction of the augmented
complex, and it is right linear; hence $\mathbb{B}$ is a projective bimodule
resolution of $T$, and ${\mathbb{B}\otimes_Ts}$ is a projective resolution of
$s$.

Let ${p\colon\mathfrak{r}^{\otimes_I3}\to T}$ be the $I$-bilinear map given by
the table of $179$ coefficients in \Cref{app:computations}, reproduced from the
appendix of~\cite{OAI26ar}: an entry $abcw$ in the row $q^d$ means that the
coefficient of the basis vector $w$ in $p(a,b,c)$ is $q^d$, where capital
letters denote dual basis vectors, so that for example ${J=j^*}$; all other
coefficients are zero. For ${\lambda\in\kk^\times}$ let $h_\lambda$ be the
automorphism ${(c,\phi)\mapsto(c,\lambda\phi)}$ of $T$, and let
${z_\lambda\colon\mathfrak{r}\to\mathfrak{r}}$ be the linear map given on basis
vectors by
\[
  z_\lambda(a)=\begin{cases}
    \lambda qa&\text{if }a\in\set{u,v,t,e^*,x^*,y^*,z^*,u^*,v^*},\\
    0&\text{otherwise}.
  \end{cases}
\]

\begin{lemma}
  \label{lemma:cocycle}
  The map $p$ has the following properties.
  \begin{enumerate}
    \item\label{it:cocycle} It is a Hochschild cocycle: for all
      ${a,b,c,d\in\mathfrak{r}}$,
      \[
        ap(b,c,d)+p(ab,c,d)+p(a,bc,d)+p(a,b,cd)+p(a,b,c)d=0,
      \]
      where $p$ is extended by zero to inputs with a factor in $I$.
    \item\label{it:weight} Every nonzero coefficient of $p(a,b,c)$ at a basis
      vector $w$, for basis vectors $a,b,c$, has exactly one more dual basis
      vector among $a,b,c$ than $w$ has.
    \item\label{it:boundary} For all ${a,b\in\mathfrak{r}}$ and
      ${\lambda\in\kk^\times}$,
      \[
        \sum_{w}p(a,b,h_\lambda(w))w^*=az_\lambda(b)+z_\lambda(ab)+
        z_\lambda(a)h_\lambda^{-1}(b),
      \]
      where $w$ runs through the radical basis vectors and $w^*$ denotes the
      dual of $w$ under the trace, so that ${(c^*)^*=c}$.
    \item\label{it:evaluation} ${p(t,x,j^*)=0}$ and ${p(t,y,j^*)=q^3f}$.
  \end{enumerate}
\end{lemma}

\begin{proof}
  These are finitely many polynomial identities in $q$ and $\lambda$, and
  they are checked by an exact computation over $\mathbb{F}_2[q,\lambda]$ on
  all inputs; see \Cref{app:computations}.
\end{proof}

Define a map ${p\colon\mathbb{B}\to\mathbb{B}[3]}$ by
\[
  p(a_0[a_1|\cdots|a_n]a_{n+1})=a_0[a_1|\cdots|a_{n-3}]\,
  p(a_{n-2},a_{n-1},a_n)a_{n+1}\quad(n\geq3),
\]
and zero on $\mathbb{B}^{-n}$ for ${n<3}$. In ${dp+pd}$, every product of adjacent
entries among the first ${n-3}$ bars appears twice, and the remaining five terms
form the left side of \Cref{lemma:cocycle}\eqref{it:cocycle}; hence $p$ is a
chain map of bimodule complexes.

\begin{proposition}
  \label{prop:cocycle-generator}
  The class $\tau$ of the chain map ${p\otimes_Ts\colon\mathbb{B}\otimes_Ts\to
  (\mathbb{B}\otimes_Ts)[3]}$ in $\operatorname{Ext}^3_T(s,s)$ is nonzero, and
  therefore generates $\operatorname{Ext}^*_T(s,s)$.
\end{proposition}

\begin{proof}
  The class is represented by the cocycle
  ${c\colon\mathbb{B}^{-3}\otimes_Ts\to s}$ obtained by composing with the
  augmentation. The element
  ${\zeta=q^2[t|x|j^*]+[t|y|j^*]}$ of ${s^{\mathrm{r}}\otimes_T\mathbb{B}^{-3}
  \otimes_Ts}$ is a cycle: its outer terms vanish because the radical
  annihilates simple modules, and since ${tx=j}$, ${ty=q^2j}$, ${xj^*=t^*}$ and
  ${yj^*=q^2t^*}$, the four inner terms cancel in pairs. A coboundary vanishes on
  every cycle, whereas ${c(\zeta)=q^2p(t,x,j^*)+p(t,y,j^*)}$ acts on $s$ as
  ${q^3\neq0}$ by \Cref{lemma:cocycle}\eqref{it:evaluation}. Hence, ${\tau\neq0}$,
  and \Cref{thm:T-ext} applies.
\end{proof}

\begin{lemma}
  \label{lemma:twist}
  For ${\lambda\in\kk^\times}$ the functor ${T_{h_\lambda}\otimes_T-}$, where
  $T_{h_\lambda}$ is $T$ with the right action twisted by $h_\lambda$, fixes
  $s$ up to a canonical isomorphism and acts on $\operatorname{Ext}^{3m}_T
  (s,s)$ as multiplication by $\lambda^{-m}$, for every ${m\geq0}$.
\end{lemma}

\begin{proof}
  The map ${T_{h_\lambda}\otimes_TM\to{}_{h_\lambda^{-1}}M}$,
  ${a\otimes m\mapsto h_\lambda^{-1}(a)m}$, is an isomorphism, natural in the
  left module $M$, onto $M$ with the action restricted along $h_\lambda^{-1}$.
  Since $h_\lambda$ is the identity modulo $DC$, which acts as zero on $s$, the
  twisted module ${}_{h_\lambda^{-1}}s$ is $s$. Applying $h_\lambda^{-1}$ to all
  entries of ${\mathbb{B}\otimes_Ts}$ defines an isomorphism from this
  resolution to its twist which induces the identity on $s$. Under it, the
  cocycle $c$ of \Cref{prop:cocycle-generator} becomes ${c\circ h_\lambda^{-1}}$.
  The only coefficients of $p$ that survive in $c$ are those at $f$ or $e$,
  which have no dual factor, and by \Cref{lemma:cocycle}\eqref{it:weight} they
  occur only on inputs with exactly one dual factor; on those inputs
  ${c\circ h_\lambda^{-1}=\lambda^{-1}c}$. Hence, the functor sends $\tau$ to
  $\lambda^{-1}\tau$, and, being exact, it preserves Yoneda products, so it
  sends $\tau^m$ to $\lambda^{-m}\tau^m$.
\end{proof}

\subsection{The tensor square}
\label{subsec:tensor-square}

Let ${E=T\otimes_\kk T}$ and ${S=s\otimes_\kk s}$, and put
${H^a}={\widehat{\operatorname{Ext}}{}^a_E(S,S)}$.

\begin{proposition}
  \label{prop:E-ext}
  The algebras $E$ and ${E^e=E\otimes_\kk E^{\op}}$ are symmetric, and
  \[
    \operatorname{Ext}^*_E(S,S)=\kk[\tau_1,\tau_2],\qquad
    \tau_1=\tau\otimes1,\ \tau_2=1\otimes\tau,
  \]
  is a polynomial algebra on two generators of degree three. Moreover, $S$ is
  not projective and ${\underline{\operatorname{End}}_E(S)=\kk}$.
\end{proposition}

\begin{proof}
  The tensor product of the symmetrising forms of the factors is a symmetric
  associative form, nondegenerate because its Gram matrix is the Kronecker
  product of two invertible matrices; the same applies to $E^e$. If
  ${\mathbb{Q}\to s}$ is a projective resolution with finitely generated terms,
  then the total complex ${\mathbb{Q}\otimes_\kk\mathbb{Q}}$ is a projective
  resolution of $S$ by finitely generated $E$-modules, as one sees by
  splitting $\mathbb{Q}$, as a complex of vector spaces, into its cohomology
  and contractible complexes. The isomorphism of complexes
  \[
    \Hom[E]{\mathbb{Q}\otimes\mathbb{Q}}{S}\cong\Hom[T]{\mathbb{Q}}{s}\otimes
    \Hom[T]{\mathbb{Q}}{s}
  \]
  gives $\operatorname{Ext}^{3m}_E(S,S)$ the basis
  $\tau_1^i\tau_2^{m-i}$, ${0\leq i\leq m}$, and all other positive degrees
  vanish. The classes $\tau_1$ and $\tau_2$ are represented by lifts of $\tau$
  acting on different tensor factors, which commute up to a sign, hence
  commute in characteristic two, so
  their monomials are the external products ${\tau^i\otimes\tau^{m-i}}$ and
  there are no relations. The last assertion follows as in \Cref{thm:T-ext}.
\end{proof}

For ${m\geq0}$ let ${P_m\subseteq\kk[\tau_1,\tau_2]}$ be the space of homogeneous
polynomials of degree $m$ in the generators, so that ${H^{3m}=P_m}$, and put
${P_m=0}$ for ${m<0}$.

\begin{proposition}
  \label{prop:E-tate}
  The Tate groups of $S$ are
  \[
    H^{3m}=P_m,\qquad H^{-3m-1}\cong DP_m\quad(m\geq0),\qquad
    H^a=0\quad\text{otherwise},
  \]
  where the isomorphism is that of \Cref{thm:tate-duality}. For ${m\geq1}$, left
  multiplication by $\tau_i$ from $H^{-3m-1}$ to $H^{-3m+2}$ is
  the transpose ${DP_m}\to{ DP_{m-1}}$ of left multiplication by $\tau_i$ from
  $P_{m-1}$ to $P_m$; it is surjective with kernel spanned by the dual of
  $\tau_{3-i}^m$. Multiplication by $\tau_i$ on ${H^{-1}=DP_0}$ is zero.
\end{proposition}

\begin{proof}
  The first assertion follows from \Cref{prop:E-ext} and Tate duality,
  \Cref{thm:tate-duality}, with ${X=Y=S}$: ${H^{-1-a}\cong DH^a}$. The
  compatibility of the duality pairing with composition~\cite[Section~2,
  (2.8)]{Lin13}, recalled after \Cref{thm:tate-duality}, identifies
  multiplication by $\tau_i$ on the negative groups with the transpose of
  multiplication by $\tau_i$ on the positive ones. Multiplication by $\tau_i$
  from $P_{m-1}$ to $P_m$ is injective with cokernel spanned by the class of
  $\tau_{3-i}^m$, which gives the description of its transpose. On $H^{-1}$
  the product lands in ${H^2=0}$.
\end{proof}

\subsection{The bimodule \texorpdfstring{$\mathcal{C}$}{C}}
\label{subsec:two-cones}

Let $\mathcal{L}$ be the mapping cone of
${p[-3]\colon\mathbb{B}[-3]\to\mathbb{B}}$,
a complex of projective $T$-bimodules with ${\mathcal{L}^j=\mathbb{B}^{j-2}\oplus
\mathbb{B}^j}$, and let
\[
  K=\mathcal{L}\otimes_\kk\mathcal{L},
\]
a complex of projective $E$-bimodules, that is, of projective $E^e$-modules,
bounded above with finitely generated terms. The complex $\mathcal{L}$ has
cohomology $T$ in degrees $0$ and $2$ and no other cohomology, and $K$ has a
filtration with four subquotients ${\mathbb{B}\otimes\mathbb{B}}$,
${(\mathbb{B}\otimes\mathbb{B})[-2]}$ twice, and ${(\mathbb{B}\otimes\mathbb{B})
[-4]}$, where ${\mathbb{B}\otimes\mathbb{B}}$ is a projective bimodule resolution
of $E$; in particular $K$ is exact in negative degrees. For ${N\geq0}$ put
\[
  C_N=\operatorname{coker}(K^{-N-1}\to K^{-N}).
\]

\begin{lemma}
  \label{lemma:two-cones}
  For ${N\geq1}$, the bimodule $C_N$ is projective as a left and as a right
  $E$-module, and ${C_{N+1}\cong C_N[-1]}$ in $\operatorname{\underline{mod}}E^e$.
  Let $\mathcal{C}$ be a finitely generated $E$-bimodule, projective on both
  sides, representing $C_N[N]$, which does not depend on $N$. Then
  ${M\mapsto\mathcal{C}M=\mathcal{C}\otimes_EM}$ induces an exact functor on
  $\operatorname{\underline{mod}}E$, and there are triangles
  \[
    S[-3]\xrightarrow{\tau_1}S\longrightarrow\mathcal{C}_1S
    \xrightarrow{\pi_1}S[-2],\qquad
    (\mathcal{C}_1S)[-3]\xrightarrow{\tau_2}\mathcal{C}_1S\longrightarrow
    \mathcal{C}S\xrightarrow{\pi_2}(\mathcal{C}_1S)[-2]
  \]
  in $\operatorname{\underline{mod}}E$, where $\mathcal{C}_1S$ is the
  evaluation at $S$ of a bimodule, projective on both sides, representing
  \[
    \operatorname{coker}\bigl((\mathcal{L}\otimes_\kk\mathbb{B})^{-2}\to
    (\mathcal{L}\otimes_\kk\mathbb{B})^{-1}\bigr)[1],
  \]
  and the second $\tau_2$ is induced by $p$ acting on the second factor.
\end{lemma}

\begin{proof}
  The cokernel functor $C_N$ is exact on degreewise split short exact
  sequences of complexes that are exact in degree ${-N-1}$: if $b$ in the
  subcomplex becomes a boundary $dv$ in the middle, the image of $v$ in the
  quotient is a cycle, hence a boundary $d\bar{w}$; replacing $v$ by
  ${v-dw}$ for a lift $w$ of $\bar{w}$ shows that $b$ is a boundary in the
  subcomplex. Hence, $C_N$ has a filtration whose factors are the cokernels
  \[
    Z_j=\operatorname{coker}\bigl((\mathbb{B}\otimes\mathbb{B})^{-j-1}\to
    (\mathbb{B}\otimes\mathbb{B})^{-j}\bigr)
  \]
  for ${j=N}$, ${N+2}$, ${N+2}$ and ${N+4}$. The
  augmented resolution ${\mathbb{B}\otimes\mathbb{B}\to E}$ splits as a complex
  of left and as a complex of right $E$-modules, since $E$ is projective on
  each side, so the $Z_j$ are projective on each side, and so is $C_N$, since
  the filtration splits on each side. Exactness of $K$ in degree ${-N}$ gives an
  exact sequence ${0\to C_{N+1}\to K^{-N}\to C_N\to0}$ with $K^{-N}$ projective
  over $E^e$, so ${C_{N+1}\cong C_N[-1]}$.

  A cosyzygy of a bimodule $M$ that is projective on each side, taken in
  $E^e$-modules, is again projective on each side: $E^e$ is symmetric, so $M$
  embeds into a projective $E^e$-module, and since modules projective on one
  side over the symmetric algebra $E$ are injective on that side, the embedding
  splits on each side. Hence, $\mathcal{C}$ can be chosen projective on both
  sides. For such a bimodule, ${\mathcal{C}\otimes_E-}$ is exact and preserves
  projective modules, and it induces an exact functor on the stable category.
  Evaluation at $S$ of the degreewise split cone sequences, which split as
  sequences of right modules, gives the two triangles; the first comes from the
  cone in the first factor of $K$, and the second from the action of $p$ on the
  second factor, as $K$ is the cone of the map induced by ${p[-3]}$ on
  ${\mathcal{L}\otimes\mathbb{B}}$.
\end{proof}

Write ${\pi=\pi_1[-2]\circ\pi_2\colon\mathcal{C}S\to S[-4]}$ for the composite
of the two projections, and put ${W^a=\widehat{\operatorname{Ext}}{}^a_E
(S,\mathcal{C}S)}$.

\begin{proposition}
  \label{prop:two-cone-profile}
  For every ${a\in\ZZ}$,
  \[
    W^a\cong\begin{cases}\kk&a\in\set{0,3},\\0&\text{otherwise,}\end{cases}
  \]
  and $\pi[3]$ induces an isomorphism ${W^3\to H^{-1}}$.
\end{proposition}

\begin{proof}
  Put ${U^a=\widehat{\operatorname{Ext}}{}^a_E(S,\mathcal{C}_1S)}$. The first
  triangle of \Cref{lemma:two-cones} gives exact sequences
  \[
    0\longrightarrow\operatorname{coker}(\tau_1\colon H^{a-3}\to H^a)
    \longrightarrow U^a\xrightarrow{\pi_1}\ker(\tau_1\colon H^{a-2}\to
    H^{a+1})\longrightarrow0.
  \]
  By \Cref{prop:E-ext,prop:E-tate}, multiplication by $\tau_1$ is injective
  on $\bigoplus_{m\geq0}P_m$ with cokernel $\kk[\tau_2]$, and on the negative
  part it is surjective with kernel spanned by the duals
  ${\varepsilon_m\in H^{-3m-1}}$ of $\tau_2^m$. Since the two possible
  contributions lie in degrees that are incongruent modulo three,
  \[
    U^{3m}=\kk x_m,\qquad U^{1-3m}=\kk y_m\quad(m\geq0),\qquad
    U^a=0\quad\text{otherwise},
  \]
  where $x_m$ is the image of $\tau_2^m$ and ${\pi_1(y_m)=\varepsilon_m}$. The
  action of $\tau_2$ on the second tensor factor commutes with all maps of the
  first triangle,
  so ${\tau_2x_m=x_{m+1}}$, ${\tau_2y_m=y_{m-1}}$ for ${m\geq1}$ and ${\tau_2y_0=0}$.
  The second triangle gives exact sequences
  \[
    0\longrightarrow\operatorname{coker}(\tau_2\colon U^{a-3}\to U^a)
    \longrightarrow W^a\xrightarrow{\pi_2}\ker(\tau_2\colon U^{a-2}\to
    U^{a+1})\longrightarrow0.
  \]
  The cokernel is nonzero only for ${a=0}$, where it is $\kk x_0$, and the kernel
  is nonzero only for ${a-2=1}$, where it is $\kk y_0$. Hence, ${W^0\cong\kk}$,
  ${W^3\cong\kk}$ via $\pi_2$ onto ${U^1=\kk y_0}$, and ${W^a=0}$ otherwise; since
  ${\pi_1(y_0)=\varepsilon_0}$ spans $H^{-1}$, the composite $\pi[3]$ is an
  isomorphism ${W^3\to H^{-1}}$.
\end{proof}

\subsection{Twisted lifts}
\label{subsec:lifts}

For ${\lambda\in\kk^\times}$ let $E_\lambda$ be the $E$-bimodule $E$ with the
right action twisted by ${h_\lambda\otimes h_\lambda}$. By \Cref{lemma:twist}
and \Cref{prop:E-ext}, the functor ${E_\lambda\otimes_E-}$ fixes $S$ and acts on
$H^{3m}$ as multiplication by $\lambda^{-m}$, for every ${m\geq0}$. Let
${\beta_0\in H^{-1}}$ be the class of the map ${S\to\Omega S}$,
${f\otimes f\mapsto f^*\otimes f^*}$, where ${\Omega S=\operatorname{rad}(E)
(f\otimes f)}$ is the kernel of the projective cover ${E(f\otimes f)\to S}$.

\begin{proposition}
  \label{prop:lifts}
  For every ${\lambda\in\kk^\times}$ there is a morphism ${g_\lambda\colon
  E_\lambda\to\mathcal{C}[3]}$ in $\operatorname{\underline{mod}}E^e$ such that
  ${\pi[3]\circ(g_\lambda\otimes_ES)=\lambda^2\beta_0}$, and ${\beta_0\neq0}$.
\end{proposition}

\begin{proof}
  Let $\mathcal{B}$ be the basis of $T$ consisting of the basis of $C$ and its
  dual basis, and let $w^*$ denote the trace dual of ${w\in\mathcal{B}}$. The
  Casimir element ${\sum_ww\otimes w^*}$ of ${T\otimes_\kk T}$ corresponds to the
  identity under ${x\otimes y\mapsto(t\mapsto x\operatorname{tr}(yt))}$, so it
  commutes with the elements of $T$. Applying $h_\lambda$ to the first factor
  and projecting to ${T\otimes_IT=\mathbb{B}^0}$ gives
  \[
    \xi_\lambda=\sum_{w\in\mathcal{B}}h_\lambda(w)\otimes_Iw^*,\qquad
    h_\lambda(a)\xi_\lambda=\xi_\lambda a\quad(a\in T),
  \]
  so ${b(x)=x\xi_\lambda}$ defines a bimodule map ${b\colon T_{h_\lambda}\to
  \mathbb{B}^0}$. Its composite with the augmentation is $\sum_wh_\lambda(w)w^*$,
  which is zero: for a basis vector $w$ of $C$, the products $ww^*$ and $w^*w$
  are the duals of the idempotents at the two ends of $w$, and each idempotent
  is the end of an even number of basis vectors of $C$ on each side, namely
  $6$ for $e$ and $4$ for $f$. Separating the same identity according to the
  left idempotent ${r\in\set{e,f}}$ of $w$ and removing the idempotent term gives
  ${\sum_{w}h_\lambda(w)w^*=r^*}$, the sum over the radical basis vectors $w$
  with left idempotent $r$; these identities are also checked in
  \Cref{app:computations}.

  Let $\mathbb{Q}$ be $\mathbb{B}$ with the right action twisted by $h_\lambda$,
  a resolution of $T_{h_\lambda}$, and define maps out of $\mathbb{Q}$ on
  generators $[a_1|\cdots|a_n]$, extended by ${\psi(a_0[\cdots]a_{n+1})}={
  a_0\psi([\cdots])h_\lambda^{-1}(a_{n+1})}$, which is the rule for bimodule
  maps from $\mathbb{Q}$. Let ${D_b\colon\mathbb{Q}\to\mathbb{B}}$ be $b$
  composed with the augmentation in degree zero and zero elsewhere, and put
  \begin{align*}
    B_n([a_1|\cdots|a_n])&=\sum_w[a_1|\cdots|a_n|h_\lambda(w)]w^*,\\
    G_n([a_1|\cdots|a_n])&=[a_1|\cdots|a_{n-1}]z_\lambda(a_n),
  \end{align*}
  with $w$ running through the radical basis vectors and ${G_0=0}$, maps of
  degrees ${-1}$ and ${+1}$. Projecting ${a\xi_\lambda=\xi_\lambda h_\lambda^{-1}(a)}$
  to the radical in the first factor gives, for a radical $a$ with right
  idempotent $r$,
  \[
    \sum_wah_\lambda(w)\otimes_Iw^*+\sum_wh_\lambda(w)\otimes_I
    w^*h_\lambda^{-1}(a)=a\otimes_Ir^*.
  \]
  Using this identity and the previous one, all terms of ${dB+Bd-D_b}$ cancel:
  the products among the first bars appear twice, the two terms involving the
  last bar of the source add up to $[a_1|\cdots|a_n]r^*$, and the product of
  the appended bar with its coefficient gives the same element. Hence,
  ${dB+Bd=0}$ in negative degrees. In degree zero, for ${r\in\set{e,f}}$,
  \[
    dB_0(r[\,]r)=\sum_wh_\lambda(w)\otimes_Iw^*+r\otimes_Ir^*=r\xi_\lambda=
    D_b(r[\,]r),
  \]
  with $w$ running through the radical basis vectors with left idempotent $r$,
  so ${dB+Bd=D_b}$. In ${dG+Gd}$ the products among the first bars cancel, and what
  remains is
  \[
    [a_1|\cdots|a_{n-2}]\bigl(a_{n-1}z_\lambda(a_n)+z_\lambda(a_{n-1}a_n)+
    z_\lambda(a_{n-1})h_\lambda^{-1}(a_n)\bigr),
  \]
  which is $pB$ by
  \Cref{lemma:cocycle}\eqref{it:boundary}; for ${n\leq1}$ both sides vanish.
  Hence, ${pB=dG+Gd}$.

  The complex ${\mathbb{Q}\otimes_\kk\mathbb{Q}}$ resolves $E_\lambda$. Define a
  map ${\Phi\colon\mathbb{Q}\otimes\mathbb{Q}\to K[3]}$ with component ${B\otimes
  D_b}$ into the summand ${(\mathbb{B}\otimes\mathbb{B})[-4][3]}$ of the underlying
  graded bimodule, component
  ${G\otimes D_b}$ into the summand corresponding to the cone in the second tensor
  factor, and zero
  components into the other two summands. The component of ${d\Phi-\Phi d}$ in
  ${(\mathbb B\otimes\mathbb B)[-4][3]}$ is ${D_b\otimes D_b}$; in the summand
  corresponding to the second tensor factor, the component ${(dG+Gd)\otimes D_b}$
  is cancelled by the term ${pB\otimes D_b}$ of the cone differential, and the
  term coming from the first factor vanishes because ${pD_b=0}$. Hence, $\Phi$ is
  a chain map in all degrees below zero, where ${D_b\otimes D_b}$ vanishes. For
  ${n\geq4}$ both complexes are exact in degree ${-n}$, so $\Phi$ induces a map
  of cokernels ${C^{-n}(\mathbb{Q}\otimes\mathbb{Q})\to C_{n-3}}$, that is, a
  stable morphism ${E_\lambda[-n]\to\mathcal{C}[3-n]}$, and shifting by $n$ gives
  $g_\lambda$. Its composite with the projection onto ${(\mathbb B\otimes\mathbb
  B)[-4][3]}$ is the
  comparison map of the bimodule map
  \[
    E_\lambda\longrightarrow\Omega_{E^e}E\subseteq(\mathbb{B}\otimes
    \mathbb{B})^0,\qquad1\longmapsto\xi_\lambda\otimes\xi_\lambda,
  \]
  which takes values in the kernel of the augmentation, since $b$ does; the
  augmentation is a projective cover, since its source is
  ${\bigoplus_rEr\otimes_\kk rE}$ over the four primitive vertex idempotents $r$
  of $E$ and the induced map on tops is an isomorphism.

  Evaluating at $S$, all terms of ${\xi_\lambda\otimes_Ts}$ vanish except the
  one with ${w^*=f}$, for which ${h_\lambda(w)=\lambda f^*}$; hence
  ${\pi[3]\circ(g_\lambda\otimes S)}$ is the class of
  ${f\otimes f\mapsto\lambda^2(f^*\otimes f^*)}$, which is $\lambda^2\beta_0$.
  Finally, ${\beta_0\neq0}$: a map ${S\to\Omega S}$ factoring through a projective
  module factors through a free module $E^r$, its components ${S\to E}$ take
  values in the left socle of $E$, and a map ${E\to\Omega S}$ is right
  multiplication by a radical element; for $u$ in the left socle and $j$ in the
  radical we have ${\operatorname{tr}(auj)=\operatorname{tr}(jau)=0}$ for all
  $a$, so ${uj=0}$, and every such factorisation is zero, whereas
  ${f^*\otimes f^*\neq0}$.
\end{proof}

\subsection{The fibre and the counterexample}
\label{subsec:fibre}

Let ${w\in W^3}$ be the element with ${\pi[3](w)=\beta_0}$, which exists and is
nonzero by \Cref{prop:two-cone-profile,prop:lifts}. For a finitely generated
$E^e$-module $M$, put
\[
  \Sigma M=\operatorname{coker}\bigl(M\to\Hom[\kk]{E^e}{M}\bigr),
\]
where $E^e$ acts on $\Hom[\kk]{E^e}{M}$ by ${(r\varphi)(t)=\varphi(tr)}$ and
${m\mapsto(t\mapsto tm)}$ is an injective homomorphism. Its target is
projective-injective, since
\[
  \Hom[\kk]{E^e}{M}\cong D(E^e)\otimes_\kk M\cong E^e\otimes_\kk M,
\]
with $M$ regarded as a vector space. If $M$ is projective on each side, so is
$\Sigma M$, as in the proof of \Cref{lemma:two-cones}, and $\Sigma M$ represents
$M[1]$. The bimodule ${\mathcal{Y}=\Sigma^4C_1}$ represents $\mathcal{C}[3]$.
Choose bimodule homomorphisms ${g_i\colon E_{H_i}\to\mathcal{Y}}$ representing
$H_i^{-2}g_{H_i}$, for ${i=1,2}$, so that both evaluate at $S$ to $w$. Let
${Q_{\mathcal{Y}}=E^e\otimes_\kk\mathcal{Y}\to\mathcal{Y}}$ be the free cover
${r\otimes y\mapsto ry}$, and define $F$ by the exact sequence
\begin{equation}
  \label{eq:fibre}
  0\longrightarrow F\longrightarrow E_{H_1}\oplus E_{H_2}\oplus Q_{\mathcal{Y}}
  \xrightarrow{(g_1,g_2,\cdot)}\mathcal{Y}\longrightarrow0,
\end{equation}
with projections ${\rho_i\colon F\to E_{H_i}}$.

\begin{proposition}
  \label{prop:fibre}
  The bimodule $F$ is finite-dimensional and projective as a left and as a
  right $E$-module. There is a unique stable morphism ${v\colon S\to FS}$ with
  ${\rho_1v=\rho_2v=\id[S]}$, and for ${V^a=\widehat{\operatorname{Ext}}{}^a_E
  (S,FS)}$ the map
  \[
    \delta^a\colon(H^a)^{\oplus2}\longrightarrow V^a,\qquad
    (g,j)\longmapsto F(g)v+v[a]j,
  \]
  is bijective for every ${a>0}$, while $\delta^0$ is surjective with kernel
  spanned by $(1,1)$.
\end{proposition}

\begin{proof}
  All terms of~\eqref{eq:fibre} are finite-dimensional, and $\mathcal{Y}$ is
  projective on each side, so the sequence splits on each side and $F$ is
  projective on each side. Tensoring with $S$ preserves exactness, the free
  bimodule evaluates to a projective module, and ${E_{H_i}\otimes_ES\cong S}$;
  hence there is a triangle
  \[
    FS\xrightarrow{(\rho_1,\rho_2)}S\oplus S\xrightarrow{(w,w)}
    \mathcal{C}S[3]\longrightarrow FS[1]
  \]
  in $\operatorname{\underline{mod}}E$, and the $\rho_i$ induce natural
  transformations ${F\otimes_E-\to E_{H_i}\otimes_E-}$ of exact functors on the
  stable category. Applying ${\widehat{\operatorname{Ext}}{}^a_E(S,-)}$ gives
  exact sequences
  \[
    (H^{a-1})^{\oplus2}\longrightarrow W^{a+2}\longrightarrow V^a
    \xrightarrow{(\rho_1,\rho_2)}(H^a)^{\oplus2}\longrightarrow W^{a+3}.
  \]
  For ${a=0}$, since ${W^2=0}$ and ${W^3=\kk w}$, the group $V^0$ maps
  isomorphically onto the kernel of ${(c,d)\mapsto(c+d)w}$, the diagonal line;
  this gives the unique $v$, which spans $V^0$. For ${a=1}$, the map
  ${(H^0)^{\oplus2}\to W^3}$ is surjective and ${W^4=0}$, and for ${a\geq2}$ both
  $W^{a+2}$ and $W^{a+3}$ vanish; hence ${(\rho_1,\rho_2)\colon V^a\to
  (H^a)^{\oplus2}}$ is an isomorphism for all ${a>0}$.

  By naturality,
  \[
    \rho_i[a]\circ F(g)\circ v=(E_{H_i}\otimes_E g)\circ\rho_iv=
    E_{H_i}\otimes_Eg\qquad\text{and}\qquad\rho_i[a]\circ v[a]\circ j=j.
  \]
  In degree zero this
  gives ${\delta^0(c,d)=(c+d)v}$, which is surjective with kernel $\kk(1,1)$. For
  ${a=3m}$ with ${m>0}$, composing $\delta^a$ with the isomorphism $(\rho_1,\rho_2)$
  gives
  \[
    (g,j)\longmapsto(H_1^{-m}g+j,\ H_2^{-m}g+j),
  \]
  by \Cref{lemma:twist,prop:E-ext}; its scalar ${2\times2}$ matrix has determinant
  ${H_1^{-m}+H_2^{-m}\neq0}$, since
  $H_1^m$ and $H_2^m$ are distinct monomials. For positive $a$ not divisible
  by three both sides vanish.
\end{proof}

\begin{theorem}[{\cite[Theorem~1.1]{OAI26ar}}]
  \label{thm:ar-counterexample}
  Let
  \[
    \Lambda=\begin{pmatrix}E&0\\F&E\end{pmatrix},
  \]
  and let $Z$ be the $\Lambda$-module of \Cref{thm:conversion} for the data of
  \Cref{prop:fibre}. Then $\Lambda$ is a finite-dimensional algebra over
  ${\kk=\mathbb{F}_2(q,H_1,H_2)}$ with eight simple modules, all one-dimensional,
  and $Z$ is a finite-dimensional, nonprojective, Gorenstein-projective
  $\Lambda$-module with ${\operatorname{Ext}^i_\Lambda(Z,Z\oplus\Lambda)=0}$ for
  all ${i\geq1}$. The same holds after extending scalars to any field
  containing~$\kk$.
\end{theorem}

\begin{proof}
  The first assertion about $Z$ is \Cref{thm:conversion} applied to
  \Cref{prop:fibre}. The radical of $E$ is ${J_E=\operatorname{rad}T\otimes T+
  T\otimes\operatorname{rad}T}$, a sum of two commuting nilpotent ideals with
  quotient $\kk^4$. The ideal of $\Lambda$ with diagonal entries in $J_E$ and
  lower left entry in $F$ is nilpotent: a product of $2n$ of its elements has
  at most one factor contributing through the lower left corner, flanked by
  ${2n-1}$ factors from $J_E$, and vanishes if ${J_E^n=0}$. Its quotient is
  $\kk^8$, so it is the radical, and $\Lambda$ has eight one-dimensional simple
  modules.

  For a field extension $K$ of $\kk$, a projective resolution of a finitely
  generated module $M$ by finitely generated projective modules gives
  \[
    K\otimes_\kk\operatorname{Ext}^i_\Lambda(M,N)\cong
    \operatorname{Ext}^i_{\Lambda_K}(M_K,N_K),
  \]
  since Hom from a finitely
  generated projective module commutes with extension of scalars. Hence, the
  vanishing of extensions persists. A complete resolution of $Z$ remains totally
  acyclic after extension of scalars, and a nonsplit projective cover sequence
  of $Z$ remains nonsplit, since its class in $\operatorname{Ext}^1$ stays
  nonzero; so $Z_K$ is Gorenstein-projective and nonprojective. The radical
  computation is unchanged.
\end{proof}

\subsection{An endomorphism algebra of infinite finitistic dimension}
\label{subsec:ar-consequences}

\begin{corollary}
  \label{coro:ar-findim}
  Let $Z'$ be an indecomposable nonprojective direct summand of the module $Z$
  of \Cref{thm:ar-counterexample}. Then the algebra
  ${\operatorname{End}_\Lambda(\Lambda\oplus Z')}$ has nine simple modules and
  infinite little finitistic dimension. The same construction gives such an
  algebra over every field containing~$\kk$.
\end{corollary}

\begin{proof}
  The module $Z'$ exists since $Z$ is nonprojective, and it inherits the
  vanishing of extensions, since Ext is additive. By \Cref{thm:ar-to-findim},
  the algebra ${\operatorname{End}_\Lambda(\Lambda\oplus Z')}$ has infinite
  little finitistic dimension. Since ${\Lambda/\operatorname{rad}\Lambda\cong
  \kk^8}$, the algebra $\Lambda$ is basic with eight indecomposable projective
  modules, and ${\Lambda\oplus Z'}$ has nine pairwise nonisomorphic
  indecomposable summands, which index the simple modules of its endomorphism
  algebra. For a field $K$ containing $\kk$, the same argument applies to
  $\Lambda_K$ and an indecomposable nonprojective summand of $Z_K$, by
  \Cref{thm:ar-counterexample}.
\end{proof}

\begin{remark}
  \label{rem:ar-sizes}
  The algebra $\Lambda$ is explicit. With the bimodule
  representatives used above, ${\dim_\kk E=400}$ and ${\dim_\kk E^e=160000}$, each
  application of $\Sigma$ multiplies dimensions by $159999$, and
  \[
    \dim_\kk F=800+159999^5\dim_\kk C_1,\qquad
    \dim_\kk\Lambda=2\dim_\kk E+\dim_\kk F,
  \]
  where ${\dim_\kk C_1=1\,623\,889\,344}$ is computed from the dimensions of the
  terms of $K$ and its cohomology (\Cref{app:computations}); hence
  ${\dim_\kk\Lambda\approx1.7\cdot10^{35}}$. A computation with minimal bimodule
  resolutions of the cones over a finite field, recorded in
  \Cref{app:computations}, determined that the two cokernels corresponding
  to $C_0$ and $C_1$ would have dimensions $784\,704$ and $1\,377\,984$, and
  was stopped before forming them, at a linear system with about
  ${3.6\cdot10^6}$ unknowns; it gives no lower bound for the dimension of a
  minimised construction. The algebra in \Cref{coro:ar-findim} has nine simple
  modules; the number ${3(l+2)}$ for \Cref{coro:main} is not determined.
\end{remark}
